% =============================================================================
%  3-Conclusiones.tex — Conclusiones y reflexión final
% =============================================================================

\section{Conclusiones}
\label{sec:conclusiones}

Este texto lo leí en mi adolescencia, más concretamente entre los años 2010 y 2011,
cuando me encontraba muy interesado en la histórica lucha de los nativos americanos en
la \textbf{\textit{Batalla de Little Bighorn}}, donde los jefes indios
\textit{Caballo Loco y Toro Sentado} lucharon contra el general Custer debido a la expansión
del viejo Oeste por parte del gobierno americano. Aunque el resultado es innegablemente una
demostración de crueldad sistemática de ambos bandos \cite{philbrick2010last},
no se puede negar que las tribus nativas americanas defendían su territorio de conquistadores
externos y que ellos conocían de primera mano las consecuencias de lo que hoy conocemos como
\textit{Deterioro Ambiental}, ya que los lagos, cañones y prácticamente cualquier entorno
natural fue arrasado o modificado por la expansión del ferrocarril, las carreteras y
posteriormente las empresas de extracción de oro y materiales de manufactura.

Hoy día, hablar de los nativos americanos es una cuestión triste: fueron relegados en su gran
mayoría a reservas o zonas naturales protegidas, lo que resulta más parecido a encerrar animales
en un zoológico que a respetar seres humanos. Los derechos de los mismos han sido aplastados
y su cultura absorbida por los medios hollywoodenses de entretenimiento visual. El territorio
no corrió con mejor suerte; aunque la carta es una muestra innegable de sabiduría y entendimiento
del desarrollo sostenible, esta fue ignorada por la clase empresarial y política que domina el mundo.
